
2.3.
Nous aurons besoin d'appliquer le foncteur @@M0@@ à des @@M1@@-Modules gradués quelconques, et non seulement à des @@M2@@-Modules gradués quasi-cohérents comme dans la construction de EGA II 3.2. Aussi allons nous indiquer rapidement comment cette construction peut s'effectuer dans le cas général.

2.3.1. Soient d'abord @@M3@@ un schéma, @@M4@@ une @@M5@@-algèbre quasi-cohérente, @@M6@@, et @@M7@@ le morphisme structural. De l'isomorphisme~:
@@M8@@
on déduit un homomorphisme canonique~:
@@M9@@
où @@M10@@ désigne l'image inverse au sens des faisceaux d'ensembles; en d'autres termes, on a un morphisme d'espaces annelés~:
@@M11@@
On en déduit un foncteur «image inverse» @@M12@@ de la catégorie des @@M13@@-Modules dans la catégorie des @@M14@@-Modules, adjoint à gauche du foncteur «image directe» @@M15@@ de la catégorie des @@M16@@-Modules dans la catégorie des @@M17@@-Modules. On a donc~:
@@M18@@
pour tout @@M19@@-Module @@M20@@.

Comme l'homomorphisme (2.3.1.1) est plat (car dans chaque fibre il correspond à une localisation), le foncteur @@M21@@ ainsi défini est un foncteur exact (cf [4], p. 164).

Si @@M22@@ est un @@M23@@-Module quasi-cohérent, cette construction donne le Module @@M24@@ défini dans EGA II 1.4.3. En effet, l'isomorphisme~:
